% ============================================================
%  USPSC-Cap3-Metodologia_3.7-Pipeline_Integrado.tex
%  Seção 3.7 — Pipeline Integrado de PLN
% ============================================================

\section{\textit{Pipeline} Integrado de PLN}
\label{sec:pipeline}

O conjunto das etapas descritas nas seções anteriores foi organizado em um \textit{pipeline} integrado de PLN de ponta a ponta (\textit{end-to-end}), implementado em Python. A modularidade da implementação garante que cada etapa possa ser executada, auditada e substituída de forma independente, o que é um requisito central para a reprodutibilidade da metodologia proposta \cite{santos2023finbert}.

\subsection{Visão Geral do \textit{Pipeline}}
\label{subsec:visao_pipeline}

O \textit{pipeline} é estruturado em seis etapas sequenciais, conforme ilustrado na \autoref{fig:pipeline}:

\begin{enumerate}
	\item \textbf{Coleta via API}: interação com o \textit{endpoint} \texttt{GET /2/users/\{id\}/tweets} da API v2 do X/Twitter, com autenticação por Bearer Token e paginação iterativa dos resultados, conforme detalhado na \autoref{subsec:coleta_api};
	
	\item \textbf{Persistência em banco de dados}: os registros retornados pela API são armazenados diretamente no banco de dados PostgreSQL — provisionado em contêiner Docker — nas tabelas
	\texttt{tweets} e \texttt{collection\_log}, eliminando a necessidade de arquivos intermediários em disco;
	
	\item \textbf{Rotulagem manual}: as publicações persistidas são exibidas sequencialmente na ferramenta de anotação (\textit{Classifier App}), onde o anotador atribui relevância financeira e polaridade de sentimento a cada registro, conforme o protocolo descrito na \autoref{sec:rotulagem};
	
	\item \textbf{AED e pré-processamento textual}: o corpus rotulado é inspecionado estatisticamente e submetido às transformações linguísticas descritas nas Seções~\ref{sec:eda} 	e~\ref{sec:preprocessamento};
	
	\item \textbf{Classificação de sentimento}: o corpus pré-processado é submetido às quatro abordagens de classificação de sentimento descritas na \autoref{sec:abordagens}: (a)~abordagem léxica via OpLexicon~v3.0; (b)~abordagem léxica via SentiLex-PT; (c)~abordagem neural via FinBERT-PT-BR; e (d)~abordagem neural via BERTimbau ajustado. Todas gravam os resultados em \texttt{tweets\_classification} com o respectivo campo \texttt{classificator}, permitindo comparação direta com o \textit{gold standard} humano na etapa de avaliação;

	\item \textbf{Agregação e análise}: as classificações individuais são agregadas temporalmente e os resultados são comparados entre as abordagens e analisados à luz do contexto macroeconômico do período, conforme apresentado no \autoref{Cap:Ava_Exp}.
\end{enumerate}

\begin{figure}[!htb]
	\centering
	\caption{\label{fig:pipeline}Diagrama do \textit{pipeline} integrado de PLN para análise de sentimento de publicações financeiras no X/Twitter}
	\resizebox{\textwidth}{!}{%
		\begin{tikzpicture}[
			node distance = 1.0cm and 0.5cm,
			box/.style = {rectangle, rounded corners=4pt, draw=black,
				fill=gray!10, text width=2.6cm, align=center,
				minimum height=1.1cm, font=\small},
			arrow/.style = {->, thick, >=stealth}
			]
	
			\node[box] (coleta) {Coleta\\via API v2};
			\node[box, right=of coleta] (db) {Persistência\\PostgreSQL};
			\node[box, right=of db] (rotul) {Rotulagem\\Manual};
			\node[box, right=of rotul] (preproc) {AED e\\Pré-proc.};
			\node[box, right=of preproc] (classif) {Classificação\\Sentimento};
			\node[box, right=of classif] (agreg) {Agregação\\e Análise};
			
			\node[above=0.35cm of coleta, font=\scriptsize] {API X/Twitter};
			\node[above=0.35cm of agreg,  font=\scriptsize] {Índice de Sentimento};
			
			\draw[arrow] (coleta)   -- (db);
			\draw[arrow] (db)       -- (rotul);
			\draw[arrow] (rotul)    -- (preproc);
			\draw[arrow] (preproc)  -- (classif);
			\draw[arrow] (classif)  -- (agreg);
			
			\node[below=0.30cm of coleta,  font=\scriptsize] {Bearer Token};
			\node[below=0.30cm of db,      font=\scriptsize] {tabela \texttt{tweets}};
			\node[below=0.30cm of rotul,   font=\scriptsize] {rótulos humanos};
			\node[below=0.30cm of preproc, font=\scriptsize] {corpus norm.};
			\node[below=0.30cm of classif, font=\scriptsize] {prob. classes};
		\end{tikzpicture}%
	}
	\legend{Fonte: Elaborado pelo autor.}
\end{figure}

\subsection{Ambiente de Execução e Ferramentas}
\label{subsec:ambiente_execucao}

A implementação foi realizada na linguagem \textbf{Python 3.11}, escolhida por sua ampla adoção na comunidade de PLN e aprendizado de máquina e pela disponibilidade de ecossistema maduro de bibliotecas especializadas. O \autoref{tab:bibliotecas} lista as principais dependências utilizadas e suas respectivas funções no \textit{pipeline}.

\begin{table}[!htb]
	\caption[Principais bibliotecas Python utilizadas no \textit{pipeline} e respectivas funções]{Principais bibliotecas Python utilizadas no \textit{pipeline} e respectivas funções}
	\label{tab:bibliotecas}
	\centering
	\begin{tabular}{ccp{10cm}}
		\hline
		\textbf{Biblioteca} & \textbf{Versão} & \textbf{Função no \textit{Pipeline}} \\
		
		\hline
		\texttt{requests} & $\geq$ 2.31 & Requisições HTTP à API v2 do X/Twitter com Bearer Token \\
		
		\hline
		\texttt{psycopg2} & $\geq$ 2.9 & Interface de comunicação com o banco de dados PostgreSQL \\
		
		\hline
		\texttt{streamlit} & $\geq$ 1.35 & Ferramenta de anotação manual (\textit{Classifier App}) e painel de AED (\texttt{eda.py}) \\
		
		\hline
		\texttt{pandas} & $\geq$ 2.0 & Manipulação tabular do corpus, AED e construção do índice temporal \\
		
		\hline
		\texttt{transformers} & $\geq$ 4.40 & Carregamento e inferência com o modelo FinBERT-PT-BR \\
		
		\hline
		\texttt{torch} & $\geq$ 2.2 & \textit{Backend} de inferência neural (PyTorch) \\
		
		\hline
		\texttt{scikit-learn} & $\geq$ 1.4 & Particionamento estratificado do corpus, cálculo de métricas de avaliação e pesos por classe para \textit{fine-tuning} \\
		
		\hline
		\texttt{accelerate} & $\geq$ 0.28 & \textit{Backend} de paralelização do \texttt{Trainer} da biblioteca \texttt{transformers}, requerido pelo protocolo de \textit{fine-tuning} do BERTimbau \\
		
		\hline
		\texttt{datasets} & $\geq$ 2.18 & Estruturação eficiente de conjuntos de dados para o ciclo de \textit{fine-tuning} supervisionado \\
		
		\hline
		\texttt{nltk} & $\geq$ 3.8 & Remoção de \textit{stopwords} para abordagem léxica \\
		
		\hline
		\texttt{spacy} & $\geq$ 3.7 & Lematização com modelo \texttt{pt\_core\_news\_lg} \\
		
		\hline
		\texttt{emoji} & $\geq$ 2.10 & Transliteração de emojis para texto \\
		
		\hline
		\texttt{plotly} & $\geq$ 5.20 & Visualizações interativas da AED e do índice de sentimento \\
		\hline
	\end{tabular}
	\legend{Fonte: Elaborado pelo autor.}
\end{table}

O código-fonte do \textit{pipeline} está organizado como um pacote Python versionado, com dependências declaradas em arquivo \texttt{requirements.txt}, de modo a facilitar a reprodução dos experimentos em ambientes distintos.